La guía pide que cada grupo diseñe una heurística propia para uno de los dos problemas
(\texttt{CornersProblem} o \texttt{FoodSearchProblem}) y compita contra las heurísticas de los
demás grupos, minimizando el número de nodos expandidos sin perder la optimalidad de la solución.

\textbf{Estado actual: pendiente de decisión del grupo.} Ya se cuenta con material completo para
competir con cualquiera de los dos problemas, con toda la documentación que exige la guía
(nombre, idea, fórmula, implementación, admisibilidad, consistencia, costo, nodos expandidos,
tiempo) ya redactada en las secciones correspondientes:

\begin{itemize}
\item \textbf{CornersProblem}: \texttt{cornersHeuristic} (Actividad 8), diámetro Manhattan de
  \{posición\} $\cup$ \{esquinas pendientes\}: 119 nodos expandidos sobre \texttt{tinyCorners},
  frente a 295 con $h(n)=0$ (factor $R=2.48$, Actividad 9).
\item \textbf{FoodSearchProblem}: \texttt{foodHeuristic} (Actividad 11, Heurística 2) --el mismo
  patrón de diámetro Manhattan, aplicado a \{posición\} $\cup$ \{comida restante\}--, 383 nodos
  expandidos sobre \texttt{testClassic} frente a 2598 de $h(n)=0$ (factor $R=6.78$).
\end{itemize}

Falta únicamente que el grupo decida cuál de las dos presentar a la competencia entre grupos; una
vez tomada esa decisión, esta sección se completa citando directamente la subsección
correspondiente (Actividad 8 o Actividad 11) como la documentación formal del reto, sin necesidad
de rehacer ningún cálculo.
